% ============================================================
% Paper deep reading: OTCA
% Source: https://arxiv.org/html/2604.19234v2
% ============================================================

\section{论文精读：OTCA}

\begin{conceptbox}
\textbf{论文：} \href{https://arxiv.org/abs/2604.19234}{Learning to Credit the Right Steps: Objective-aware Process Optimization for Visual Generation}\\
\textbf{简称：} OTCA。\\
\textbf{版本：} arXiv:2604.19234v2，2026-04-27。\\
\textbf{方向标签：} Flow Matching、Visual GRPO、Credit Assignment、Multi-objective Reward、Image / Video Generation。\\
\textbf{阅读日期：} 2026-04-28。\\
\textbf{阅读口径：} 只放在我们当前的主线里读：\textbf{连续 flow / visual GRPO 的信用分配}。论文覆盖 image 和 video，术语里会说 denoising trajectory；这里统一理解成 Flow-SDE 采样轨迹上的 policy optimization。
\end{conceptbox}

\begin{summarybox}
\textbf{一句话总结：} OTCA 把原本一个 trajectory-level advantage 拆成两个问题：\textbf{什么时候更新}和\textbf{更新哪个 reward objective}。核心公式是
\[
\tilde{A}_{t}^{i}
=
w_t^i
\sum_{k=1}^{K} c_k^i A_k^i .
\]
其中 \(w_t^i\) 来自 TCD，负责 temporal credit；\(c_k^i\) 来自 MOCA，负责多目标 reward credit。它和我们前面几篇的关系是：DenseGRPO / TempFlow 更偏 reward gain 或 branch 反事实，PCPO 更偏 ratio 侧权重，OTCA 则进一步问 multi-reward setting 下 advantage 本身应该怎么组织。
\end{summarybox}

\begin{center}
  \includegraphics[width=0.95\linewidth]{otca/fig1_framework_comparison.png}

  {\small \textbf{图 1：} OTCA 和普通 GRPO 的区别。普通做法把多个 reward 压成一个 scalar，再广播给所有 timesteps；OTCA 试图同时做 timestep credit 和 objective credit。图源：arXiv HTML。}
\end{center}

\subsection{0. 它接在 PCPO 后面读的原因}

我们这条主线一直在拆同一个问题：

\[
\text{terminal reward 太粗，不能说明哪个 timestep 真正导致好结果。}
\]

前四篇可以这样放：

\begin{center}
\resizebox{\linewidth}{!}{
\begin{tabular}{lll}
\toprule
\textbf{论文} & \textbf{主要改动的位置} & \textbf{它问的问题}\\
\midrule
DenseGRPO & reward gain & 哪个 step 真的提升了最终 reward？\\
Chunk-GRPO & optimization unit & step 太细、trajectory 太粗时，能不能按 chunk 更新？\\
TempFlow-GRPO & branch / counterfactual & 某个时间点的分叉到底带来什么影响？\\
PCPO & policy ratio / objective scale & ratio 数学形式是否隐式偏向某些 steps？\\
OTCA & temporal credit + reward objective & 哪个 step 重要，以及这个 step 应该优化哪个目标？\\
\bottomrule
\end{tabular}
}
\end{center}

所以 OTCA 不是在重复 DenseGRPO。它的切入点是 multi-reward visual generation：

\[
\begin{aligned}
R
&=
\text{CLIP / text alignment}
+
\text{aesthetic / visual quality}\\
&\quad
+
\text{motion / temporal consistency}
+
\cdots
\end{aligned}
\]

如果简单把它们加成一个 scalar，再把这个 scalar advantage 平均给所有 timesteps，就会同时犯两个错误：

\begin{enumerate}[leftmargin=1.5em]
  \item 不知道 \textbf{哪个 timestep} 应该承担更大 credit。
  \item 不知道 \textbf{哪个 reward objective} 在当前生成阶段更应该被强调。
\end{enumerate}

\begin{intubox}
\textbf{直觉：} 早期 denoising 可能更影响 layout / semantic structure，中后期可能更影响 texture / aesthetic / motion details。把所有目标都压成一个 \(A_i\)，再平等广播给所有 steps，其实是在假设“每个阶段都应该优化同一种东西”，这个假设很粗。
\end{intubox}

\subsection{1. 标准 GRPO 的问题：一个 advantage 管所有 steps}

普通 visual GRPO 对同一个 prompt 采样 \(G\) 个样本，得到每个样本的 reward：

\[
r_1,r_2,\ldots,r_G .
\]

然后做 group-wise normalized advantage：

\[
A_i
=
\frac{
r_i-\mathrm{mean}(r_1,r_2,\ldots,r_G)
}{
\mathrm{std}(r_1,r_2,\ldots,r_G)
}.
\]

训练目标是：

\[
\mathcal{J}(\theta)
=
\mathbb{E}
\left[
\frac{1}{G}\sum_{i=1}^{G}
\frac{1}{T}\sum_{t=1}^{T}
\min
\left(
\rho_t^i A_i,\,
\mathrm{clip}(\rho_t^i,1-\epsilon,1+\epsilon)A_i
\right)
\right].
\]

这里的问题非常直接：

\[
A_i
\quad
\text{和}
\quad
t
\quad
\text{无关。}
\]

也就是说，一张图最后 reward 高，整条 trajectory 所有 steps 都被正向更新；一张图最后 reward 低，整条 trajectory 所有 steps 都被压低。OTCA 要把这个单一的 \(A_i\) 换成 timestep-specific effective advantage：

\[
A_i
\quad \Longrightarrow \quad
\tilde{A}_t^i .
\]

\subsection{2. OTCA 总公式：when 乘 what}

OTCA 的核心公式是：

\[
\tilde{A}_{t}^{i}
=
w_t^i
\cdot
\sum_{k=1}^{K} c_k^i A_k^i .
\]

这个式子可以直接读成：

\[
\text{effective advantage}
=
\text{temporal importance}
\times
\text{objective mixture advantage}.
\]

各个符号含义是：

\begin{itemize}[leftmargin=1.5em]
  \item \(i\)：同一个 prompt 下的第 \(i\) 条 sampled trajectory。
  \item \(t\)：flow / denoising trajectory 上的某个 timestep。
  \item \(k\)：第 \(k\) 个 reward objective，例如 CLIP-T、Aesthetic、HPS，或者 video 里的 visual / motion / text alignment。
  \item \(A_k^i\)：第 \(i\) 个样本在第 \(k\) 个 reward objective 上的 group-normalized advantage。
  \item \(w_t^i\)：第 \(t\) 个 step 对最终结果的相对贡献。
  \item \(c_k^i\)：第 \(i\) 条 trajectory 上，不同 reward objectives 的融合权重。
\end{itemize}

\begin{center}
  \includegraphics[width=0.92\linewidth]{otca/fig2_method_overview.png}

  {\small \textbf{图 2：} OTCA 方法概览。TCD 产生 timestep 权重 \(w_t^i\)，MOCA 产生多目标权重 \(c_k^i\)，最后合成 \(\tilde{A}_t^i\)。图源：arXiv HTML。}
\end{center}

\begin{warnbox}
\textbf{一个细节要注意：} 论文文字里常说 timestep-aware objective mixture，但主公式里的 MOCA 权重是 \(c_k^i\)，不是 \(c_{t,k}^i\)。也就是说，严格从公式看，OTCA 的 objective mixture 是 sample / trajectory conditioned；真正随 timestep 变化的是 \(w_t^i\)。它不是完全为每个 timestep 单独求一组 objective weights。
\end{warnbox}

\subsection{3. Flow-SDE 背景：为什么有 trajectory 可以做 credit}

Flow matching 的确定性采样通常写成 ODE：

\[
\mathrm{d}\mathbf{z}_t
=
\mathbf{u}_t\,\mathrm{d}t .
\]

但 GRPO 需要 stochastic trajectory，才能把每一步看成 policy transition 并计算 ratio。因此论文沿用 prior work，把 ODE 换成 reverse-time SDE：

\[
\mathrm{d}\mathbf{z}_t
=
\left(
\mathbf{u}_t
-
\frac{1}{2}\varepsilon_t^2
\nabla_{\mathbf{z}_t}\log p_t(\mathbf{z}_t)
\right)\mathrm{d}t
+
\varepsilon_t\,\mathrm{d}\mathbf{w}.
\]

这样每个 step 都可以被理解为：

\[
\mathbf{a}_{t,i}
\sim
\pi_\theta(\cdot|\mathbf{s}_{t,i}),
\]

于是可以写：

\[
\rho_t^i
=
\frac{
\pi_\theta(\mathbf{a}_{t,i}|\mathbf{s}_{t,i})
}{
\pi_{\theta_{\mathrm{old}}}(\mathbf{a}_{t,i}|\mathbf{s}_{t,i})
}.
\]

\textbf{这一点和 PCPO 是同一个入口：} 都先把 flow ODE 放到 SDE / stochastic policy 的框架里，再谈 GRPO objective。区别是 PCPO 研究 \(\rho_t^i\) 本身的 hidden weight；OTCA 研究 \(\tilde{A}_t^i\) 应该怎么构造。

\subsection{4. TCD：用 latent similarity 估计哪个 step 重要}

TCD 全称是 Trajectory-Level Credit Decomposition。它负责估计：

\[
w_t^i
=
\text{第 \(i\) 条 trajectory 的第 \(t\) 个 step 应该拿多少 temporal credit。}
\]

它没有像 DenseGRPO 那样对每个 \(x_t\) 做 ODE readout 再打 reward，而是用一个更便宜的 latent proxy。

\subsubsection{4.1 先看中间 latent 和 final latent 的相似度}

对第 \(i\) 条 trajectory，记中间 latent 为 \(z_t^i\)，最终 latent 为 \(z_{\mathrm{final}}^i\)。TCD 计算它们的 cosine similarity：

\[
S_t^i
=
\frac{
z_t^i \cdot z_{\mathrm{final}}^i
}{
\|z_t^i\|\,\|z_{\mathrm{final}}^i\|
}.
\]

这个 \(S_t^i\) 的直觉是：

\[
\text{当前 latent 和最终结果方向越接近，说明它越像最终要生成的东西。}
\]

但论文不直接用 \(S_t^i\)，而是用相邻两步的变化量：

\[
\Delta S_t^i
=
S_t^i-S_{t+1}^i .
\]

注意这里的符号和时间方向有关。因为 denoising / flow trajectory 通常从 noisy 走向 clean，论文按它自己的 \(t\) 编号写成 \(S_t^i-S_{t+1}^i\)。我们读的时候只要抓住意思：

\[
\Delta S_t^i
\approx
\text{这个 transition 让 latent 朝 final latent 靠近了多少。}
\]

如果某一步让 latent 更明显地对齐最终状态，就认为它对最终结果更重要。

\subsubsection{4.2 把 similarity gain 变成 temporal weight}

OTCA 把负的贡献截掉，然后归一化：

\[
w_t^i
=
\frac{
\max(0,\Delta S_t^i)+\epsilon
}{
\sum_{\tau}
\left(
\max(0,\Delta S_\tau^i)+\epsilon
\right)
}.
\]

这保证：

\[
w_t^i \ge 0,
\qquad
\sum_t w_t^i = 1.
\]

于是每条 trajectory 的 temporal credit 被重新分配到不同 timesteps：

\[
\text{重要 steps}
\Rightarrow
\text{更大的 }w_t^i,
\qquad
\text{不重要 steps}
\Rightarrow
\text{更小的 }w_t^i.
\]

\begin{intubox}
\textbf{和 DenseGRPO 的区别：} DenseGRPO 的 \(\Delta R_t\) 更像真实 reward gain，但要 ODE readout 和 reward scoring，成本高。OTCA 的 \(\Delta S_t\) 是 latent proxy，不需要额外 reward model 调每个中间点，便宜很多，但因果性弱一些。
\end{intubox}

\begin{center}
  \includegraphics[width=0.9\linewidth]{otca/fig6_tcd_proxy_validation.png}

  {\small \textbf{图 3：} 论文在 appendix 里用 \(\Delta S_t^i\) 和真实 reward gain \(\Delta R_t^i\) 的趋势对比来支持 TCD。它们大体相关，但不是严格相等。图源：arXiv HTML。}
\end{center}

\subsection{5. MOCA：多目标 reward 不要粗暴相加}

MOCA 全称是 Multi-Objective Credit Allocation。它负责估计：

\[
c_k^i
=
\text{第 \(i\) 条 trajectory 上第 \(k\) 个 reward objective 的权重。}
\]

假设我们有 \(K\) 个 reward objectives。对第 \(i\) 个样本，每个 objective 都可以得到一个 group-normalized advantage：

\[
\mathbf{A}^i
=
\left[
A_1^i,A_2^i,\ldots,A_K^i
\right].
\]

如果直接加权平均：

\[
A_{\mathrm{sum}}^i
=
\sum_{k=1}^{K}\alpha_k A_k^i,
\]

就有两个问题：

\begin{enumerate}[leftmargin=1.5em]
  \item 不同 reward 可能冲突，比如 aesthetic 上升但 text alignment 下降。
  \item 固定 \(\alpha_k\) 不知道当前样本和当前 trajectory 更需要哪个 objective。
\end{enumerate}

\subsubsection{5.1 从 gradient-space MOO 到 advantage-space MOO}

标准 multi-objective optimization 可以在 gradient space 里做：

\[
\min_{\{c_k\}_{k=1}^{K}}
\left\|
\sum_{k=1}^{K}
c_k
\nabla_\theta \mathcal{L}_k(\theta)
\right\|^2,
\qquad
\sum_{k=1}^{K}c_k=1,\quad c_k\ge 0.
\]

但对大视觉生成模型来说，为每个 objective 单独算 gradient 再解 MOO 太贵。OTCA 的观察是：对同一个 sample，在同一次 policy update 里，不同 objective 的梯度主要差在 scalar advantage \(A_k^i\)。于是把问题投影到 advantage space：

\[
\min_{\{c_k^i\}_{k=1}^{K}}
\left\|
\sum_{k=1}^{K}
c_k^i A_k^i
\right\|^2,
\qquad
\sum_{k=1}^{K}c_k^i=1,\quad c_k^i\ge 0.
\]

这就是 MOCA 的主体：不用真正存每个 objective 的梯度，而是在低维 advantage vector 上决定 reward mixture。

\subsubsection{5.2 为什么还要 exploration bias}

单纯 minimum-norm 有一个问题：它容易太保守。因为它倾向于找到一个让综合 advantage 接近 0 的组合，避免冲突，但也可能把真正高收益的方向抵消掉。

所以 OTCA 构造了一个 trajectory-level exploration signal：

\[
q_i
=
\sum_{t=1}^{T} w_t^i \Delta S_t^i,
\qquad
e_i
=
\frac{
|q_i|
}{
\mathrm{std}(\{q_j\}_{j=1}^{G})+\epsilon
}.
\]

其中 \(q_i\) 表示这条 trajectory 的结构变化强度，\(e_i\) 表示它相对组内其他 samples 有多 distinctive。然后定义 exploration strength：

\[
\lambda_i
=
\frac{
\max\left(0,\sum_{k=1}^{K} A_k^i\right)
}{
\left|\sum_{k=1}^{K} A_k^i\right|+\epsilon
}
\cdot
\tanh(e_i).
\]

这个设计有两个门控：

\begin{itemize}[leftmargin=1.5em]
  \item 只有综合 advantage 为正时才鼓励探索，避免放大坏方向。
  \item \(\tanh(e_i)\) 把探索强度压到有界范围，避免异常 trajectory 造成不稳定。
\end{itemize}

最终的 MOCA solver 变成：

\[
\min_{\{c_k^i\}}
\left\|
\sum_{k=1}^{K} c_k^i A_k^i
\right\|^2
-
\lambda_i
\sum_{k=1}^{K} c_k^i A_k^i,
\qquad
\sum_{k=1}^{K}c_k^i=1,\quad c_k^i\ge 0.
\]

前半项负责 conflict-aware fusion，后半项负责不要太保守。

\subsection{6. 最终训练目标}

OTCA 最后只是把原本的 \(A_i\) 换成 \(\tilde{A}_t^i\)：

\[
\mathcal{J}_{\mathrm{OTCA}}(\theta)
=
\mathbb{E}
\left[
\frac{1}{G}\sum_{i=1}^{G}
\frac{1}{T}\sum_{t=1}^{T}
\min
\left(
\rho_t^i \tilde{A}_t^i,\,
\mathrm{clip}(\rho_t^i,1-\epsilon,1+\epsilon)\tilde{A}_t^i
\right)
\right].
\]

把 \(\tilde{A}_t^i\) 展开就是：

\[
\rho_t^i
\left(
w_t^i
\sum_{k=1}^{K}c_k^i A_k^i
\right).
\]

所以每个 step 的实际更新强度由三部分共同决定：

\[
\text{policy ratio}
\times
\text{temporal credit}
\times
\text{objective credit}.
\]

\begin{summarybox}
\textbf{和 PCPO 的连接：} PCPO 提醒我们 \(\rho_t^i\) 里也有 timestep bias；OTCA 处理的是 \(\tilde{A}_t^i\) 的结构。更完整的 credit assignment 其实需要同时看 ratio 侧和 advantage 侧。
\end{summarybox}

\subsection{7. 实验怎么证明它有效}

\subsubsection{7.1 Image generation}

图像实验用 FLUX.1-dev。训练 reward 包括 CLIP-T、LAION-Aesthetic、HPSv2.1；同时用 PickScore 和 ImageReward 看泛化。

\begin{center}
\resizebox{\linewidth}{!}{
\begin{tabular}{lccccc}
\toprule
\textbf{Method} & \textbf{CLIP-T} & \textbf{Aesthetic} & \textbf{HPS} & \textbf{PickScore} & \textbf{ImageReward}\\
\midrule
Flux & 0.2682 & 6.2508 & 0.2986 & 22.66 & 1.1001\\
DanceGRPO & 0.2619 & \textbf{6.8917} & 0.3018 & 22.70 & 1.0172\\
VIPO & 0.2722 & 6.7813 & 0.3026 & 22.69 & 1.1128\\
OTCA & \textbf{0.3071} & 6.6028 & \textbf{0.3225} & \textbf{22.97} & \textbf{1.1998}\\
\bottomrule
\end{tabular}
}

{\small \textbf{表 1：} 图像指标。OTCA 在 CLIP-T、HPS、PickScore、ImageReward 上最好；Aesthetic 是 DanceGRPO 更高。}
\end{center}

这里最值得注意的不是“每个指标都赢”，而是 \textbf{trade-off 更均衡}。DanceGRPO 的 aesthetic 很高，但 semantic / preference 指标不一定一起变好；OTCA 通过 MOCA 把多个 objective 更平衡地组织起来。

\begin{center}
  \includegraphics[width=0.95\linewidth]{otca/fig3_image_qualitative.png}

  {\small \textbf{图 4：} 图像生成质性结果。OTCA 主要展示语义一致性、结构清晰度和局部细节上的提升。图源：arXiv HTML。}
\end{center}

\subsubsection{7.2 Video generation}

视频实验用 Wan2.2-T2V-14B-480P，reward 来自 VideoAlign，包含 visual quality、motion quality、text alignment。

\begin{center}
\resizebox{\linewidth}{!}{
\begin{tabular}{lccccccc}
\toprule
\textbf{Method} & \textbf{Color} & \textbf{Dynamic} & \textbf{Imaging} & \textbf{Multiple} & \textbf{Scene} & \textbf{Spatial} & \textbf{Subject}\\
\midrule
Wan2.2 & 89.79 & 70.83 & 68.34 & 63.79 & 35.90 & 80.54 & 93.78\\
DanceGRPO & 90.35 & 71.45 & 67.93 & 64.85 & 35.18 & 83.25 & 93.76\\
OTCA & \textbf{92.88} & \textbf{75.00} & \textbf{68.58} & \textbf{68.45} & \textbf{38.45} & \textbf{88.57} & \textbf{93.95}\\
\bottomrule
\end{tabular}
}

{\small \textbf{表 2：} VBench 细分指标。OTCA 在列出的 video 维度上都优于 Wan2.2 和 DanceGRPO。}
\end{center}

\begin{center}
\resizebox{0.7\linewidth}{!}{
\begin{tabular}{lccc}
\toprule
\textbf{Method} & \textbf{Quality} & \textbf{Semantic} & \textbf{Total}\\
\midrule
Wan2.2 & 83.30 & 73.10 & 81.26\\
DanceGRPO & 83.31 & 73.31 & 81.31\\
OTCA & \textbf{83.72} & \textbf{75.19} & \textbf{82.01}\\
\bottomrule
\end{tabular}
}

{\small \textbf{表 3：} VBench 总分。提升主要体现在 semantic score。}
\end{center}

\begin{center}
  \includegraphics[width=0.95\linewidth]{otca/fig4_video_qualitative.png}

  {\small \textbf{图 5：} 视频生成质性结果。论文强调 OTCA 改善物理合理性、运动演化和语义一致性。图源：arXiv HTML。}
\end{center}

\subsubsection{7.3 Ablation：TCD 和 MOCA 都有用}

\begin{center}
\resizebox{\linewidth}{!}{
\begin{tabular}{lccccc}
\toprule
\textbf{Ablation} & \textbf{CLIP-T} & \textbf{Aesthetic} & \textbf{HPS} & \textbf{PickScore} & \textbf{ImageReward}\\
\midrule
Flux & 0.2682 & 6.2508 & 0.2986 & 22.66 & 1.1001\\
MOCA-only & 0.2787 & 6.2120 & 0.3049 & 22.64 & 1.1250\\
TCD-only & 0.2708 & 6.3270 & 0.3104 & 22.87 & \textbf{1.2618}\\
OTCA Full & \textbf{0.3071} & \textbf{6.6028} & \textbf{0.3225} & \textbf{22.97} & 1.1998\\
\bottomrule
\end{tabular}
}

{\small \textbf{表 4：} Ablation。TCD-only 和 MOCA-only 都有提升，完整 OTCA 在多数指标上最好。}
\end{center}

\begin{center}
  \includegraphics[width=0.95\linewidth]{otca/fig5_reward_curves.png}

  {\small \textbf{图 6：} Reward curves。完整 OTCA 的收敛速度和最终 reward 普遍更好。图源：arXiv HTML。}
\end{center}

\subsection{8. 训练成本和实现细节}

论文给出的主要实现设置是：

\begin{itemize}[leftmargin=1.5em]
  \item \textbf{Image：} FLUX.1-dev，group size \(G=12\)，训练分辨率 \(512\times512\)，16 sampling steps。
  \item \textbf{Video：} Wan2.2-T2V-14B-480P，group size \(G=32\)，训练分辨率 \(240\times416\times53\)，16 sampling steps。
  \item \textbf{Inference：} Flux 用 \(512\times512\)、50 steps；Wan2.2 用 \(480\times832\times53\)、50 steps。
  \item \textbf{Hardware：} image experiments 用 8 张 NVIDIA H100；video experiments 用 32 张 NVIDIA H100。
  \item \textbf{SDE stochasticity：} \(\eta=0.3\) for Flux，\(\eta=0.8\) for Wan2.2。
  \item \textbf{Learning rate：} image \(1\times10^{-5}\)，video \(5\times10^{-6}\)。
\end{itemize}

\begin{warnbox}
\textbf{成本判断：} OTCA 比 DenseGRPO 那种每个 timestep ODE readout 便宜，因为 TCD 用 latent similarity proxy；但它并不是低成本小实验，尤其视频部分仍然是 32 张 H100 级别。它的主要节省来自“不对每个中间 latent 都 decode / reward scoring”，不是来自 GRPO 本身变便宜。
\end{warnbox}

\subsection{9. 我对这篇的理解和保留意见}

\subsubsection{9.1 最值得学的点}

\textbf{OTCA 的最大价值是把 credit assignment 拆成了两个正交轴：}

\[
\text{when}
\quad
\text{and}
\quad
\text{what}.
\]

这比只说“每个 step 都有不同 reward”更进一步。因为在真实 visual post-training 里，reward 往往不是一个，而是一组：

\[
\text{text alignment},\quad
\text{aesthetic},\quad
\text{human preference},\quad
\text{motion},\quad
\text{physical plausibility}.
\]

如果只做 temporal credit，不处理 reward objective 冲突，仍然可能把某些 steps 往错误目标上推。

\subsubsection{9.2 最大的局限}

\begin{enumerate}[leftmargin=1.5em]
  \item \textbf{TCD 是 proxy，不是真正 causal reward gain。} \(\Delta S_t^i\) 衡量的是 latent 朝 final latent 的对齐变化，不等价于 reward improvement。它可能和最终 reward 有相关性，但不是 DenseGRPO 式的 readout counterfactual。
  \item \textbf{MOCA 的 \(c_k^i\) 不是逐 timestep 的。} 主公式里 \(c_k^i\) 没有 \(t\)，所以它更多是在 trajectory / sample 层面做 objective mixture，再通过 \(w_t^i\) 分配到各 steps。
  \item \textbf{advantage-space MOO 是近似。} 它避免了 per-objective gradient 的计算，但也意味着它没有真正解决 gradient conflict，只是在 scalar advantage 层面做了一个 proxy。
  \item \textbf{实验成本仍然很高。} 特别是视频部分用了 32 张 H100，这说明方法虽然避免了 Dense readout，但 visual GRPO 本身仍是重计算路线。
\end{enumerate}

\subsection{10. 和我们当前主线的关系}

读完 OTCA 后，这条主线可以更新成：

\begin{center}
\resizebox{\linewidth}{!}{
\begin{tabular}{llll}
\toprule
\textbf{层面} & \textbf{代表论文} & \textbf{解决什么} & \textbf{还没解决什么}\\
\midrule
step reward gain & DenseGRPO & 每步真实 reward improvement & 成本高，readout 太密\\
chunk unit & Chunk-GRPO & 更新粒度更平滑 & 仍多用 terminal advantage\\
branch counterfactual & TempFlow-GRPO & 单步分叉的局部影响 & branch 设计和成本复杂\\
ratio / objective scale & PCPO & ratio 自带 timestep bias & 不直接重构 reward credit\\
temporal + objective credit & OTCA & when + what & TCD / MOCA 都是 proxy\\
\bottomrule
\end{tabular}
}
\end{center}

\begin{summarybox}
\textbf{当前阶段我的判断：} OTCA 值得学，但不要把它看成 credit assignment 的最终答案。它更像一个框架提醒：真实 visual GRPO 的有效 advantage 应该同时包含 temporal importance 和 objective importance。对于我们后续想做的改进，更关键的问题仍然是：能否用低成本方式估计更接近 causal reward gain 的 \(w_t\) 或 chunk-level \(w_{\mathrm{ch}}\)，同时保留多目标 reward 的动态分配。
\end{summarybox}
